\section{Existing Solutions}
\label{sec:existing_solutions}

\subsection{Evolution of Sequential Recommendation Models}

Sequential recommendation has evolved through several model families, each addressing a limitation of its predecessor.

\paragraph{GRU4Rec —}
Hidasi et al.~\cite{hidasi2016gru4rec} first demonstrated that recurrent neural networks could outperform matrix factorisation on session-based recommendation tasks.
A gated recurrent unit (GRU) encoder reads the interaction sequence and produces a hidden state that is used to rank candidate items.
The principal limitation is the sequential nature of the GRU: the hidden state at each step depends on the previous step, preventing parallelisation during training and creating an information bottleneck for long sequences.

\paragraph{SASRec —}
Kang and McAuley~\cite{kang2018self} replaced the GRU with a causal self-attention Transformer, enabling parallel computation over the full sequence during training.
SASRec significantly outperforms GRU-based models on standard benchmarks and is substantially faster to train.
However, it uses only item ID embeddings and does not incorporate item content features, limiting its ability to generalise to new or rare items.

\paragraph{BERT4Rec —}
Sun et al.~\cite{sun2019bert4rec} applied bidirectional attention with a Cloze-style masked-item objective.
Access to future context improves accuracy on held-out items but requires a slower, more complex inference procedure (masking the target position and re-encoding).

\paragraph{DIF-SASRec —}
Xie et al.~\cite{xie2022difsasrec} observed that prior methods fuse side information (text, category) into item embeddings before self-attention, which prevents the model from differentiating between content-relevant and co-occurrence-relevant attention patterns.
DIF-SASRec decouples these streams into separate attention modules and shows consistent gains over SASRec and BERT4Rec across multiple datasets.
This architecture is adopted as Pipeline~B of the proposed system.

\subsection{The Prior System: GRU-SeqDQN}

The system evaluated in this report supersedes a prior implementation (hereafter GRU-SeqDQN) that paired a GRU sequence encoder with a Deep Q-Network (DQN) policy for recommendation.
The DQN formulates recommendation as a reinforcement learning problem: the agent receives a reward signal from user interactions and learns a Q-function mapping (user state, candidate item) pairs to expected return.

While this formulation is theoretically appealing, the evaluation reported in Chapter~5 shows that GRU-SeqDQN achieves an HR@10 of only 0.1031 — near the random baseline of 0.1000 — on the 100{,}000-user evaluation set.
Two structural limitations explain this failure.
First, the GRU encoder inherits the sequential bottleneck identified above.
Second, DQN training is highly sample-inefficient at the scale of millions of items: the action space is too large for the Q-function to converge reliably in a reasonable number of training steps.
The proposed DIF-SASRec (Pipeline B) and Cleora-based (Pipeline A) pipelines directly address both limitations.

